% =====================================================================
%  Chapter 10 -- Conclusion
% =====================================================================
\chapter{Conclusion and Future Scope}\label{ch:conclusion}\label{ch:future}

We asked whether a simple electronic circuit can show that the speed of a change, not only its size,
decides the final state. For a Wien-bridge oscillator whose gain is swept across threshold, it can.

Three results stand out. The switching lag grows as $\sqrt{\lambda}$, with a switch-off coefficient that
matches a parameter-free prediction in simulation and on the bench. In simulation, fast sweeps of the
feedback resistor change the long-time state from bursting to permanent oscillation. On the bench,
sweeps of the input resistor along one fixed path show the signature of rate-induced tipping: below about
\SI{6.1}{\kilo\ohm\per\second} the oscillation re-ignites every cycle, above it the circuit falls silent
for good, the two separated by a $2\%$ window in rate. All of these follow from the growth budget of
\cref{sec:one-idea}.

The Wien-bridge oscillator is thus a controllable, quantitatively understood laboratory system showing
clear signatures of rate-induced tipping.
